% !TeX root = ../../金融经济学_王江_讲义及习题.tex
% 10_完全市场 课后习题（由章节文件拆分）

\par\noindent\begin{minipage}{\linewidth}
\begin{problem}
\WJTenExercise{exer:wj-ten-1}{*} 考虑一个两期（$t=0,\,1$）、两个状态（$\omega=a,\,b$），存在 $K$ 个参与者的经济。状态 $a$ 发生的概率为 $\pi_{a}$ 而状态 $b$ 发生的概率为 $\pi_{b}$。所有参与者都是不满足的且严格风险厌恶的。他们期望效用有如下形式：
\[
    u_{k}(c_{k,0})+\pi_{a}u_{k}(c_{k,a})+\pi_{b}u_{k}(c_{k,b}),\quad k=1,\dots,K.
\]

参与者 $k$ 的禀赋是 $\begin{bmatrix}e_{k,0} \\ e_{k,a} \\ e_{k,b}\end{bmatrix}$。整个经济的总禀赋是 $\begin{bmatrix}C_{0} \\ C_{a} \\ C_{b}\end{bmatrix}$，其中 $C_{s}=\sum_{k}e_{k,s}$，$s=0,a,b$。假设 $C_{a}>C_{b}$。假设市场上交易的 Arrow--Debreu 状态或有要求权的集合是完全的。令 $\phi_{a}$ 和 $\phi_{b}$ 分别是状态 $a$ 和状态 $b$ 的状态价格。
\begin{enumerate}[label=(\alph*)]
    \item 定义这个经济的均衡。简要讨论均衡的存在性和唯一性。
    \item 令代表性参与者的（间接）期望效用函数为
    \[
        u(C_{0})+\pi_{a}u(C_{a})+\pi_{b}u(C_{b})
    \]
    且 $u'''(\cdot)<0$。当以下变量变化而其它变量保持不变时，状态价格（$\phi_{a}$ 和 $\phi_{b}$）、它们的比例（$\phi_{a}/\phi_{b}$）和无风险利率如何变化（上升或下降）？给出证明。
    \begin{itemize}
        \item $C_{0}$ 上升；
        \item $C_{a}$ 上升；
        \item $C_{b}$ 上升；
        \item $C_{a}$ 和 $C_{b}$ 同比例上升；
        \item $C_{a}-C_{b}$ 上升（禀赋的风险上升）。
    \end{itemize}
\end{enumerate}
比较禀赋风险增加时，保持 $\pi_aC_a+\pi_bC_b$ 不变。
\end{problem}
\end{minipage}\par
\Needspace{5\baselineskip}
\begin{solution}
\begin{enumerate}[label=(\alph*)]
 \item 均衡由正状态价格 $(\phi_a,\phi_b)$ 与各人消费组成。给定价格，各人最大化原题效用，满足
 \[
 c_{k,0}+\phi_ac_{k,a}+\phi_bc_{k,b}
 \le e_{k,0}+\phi_ae_{k,a}+\phi_be_{k,b},
 \]

 并且 $\sum_kc_{k,s}=C_s$（$s=0,a,b$）。在偏好连续、凸、严格单调及个人禀赋严格为正等均衡存在条件下，完全市场存在均衡。严格凹性使给定价格下的个人选择唯一；一般均衡的唯一性还取决于总超额需求。\ifdefined\WJHandoutKnowledge 正文定理~\ref{thm:wj-ten-3} 后给出了多重均衡的例子。\else 完整讲义的均衡存在性讨论中给出了多重均衡的例子。\fi

 \item 以题设代表性效用 $u$ 保持不变为前提，
 \[
 \phi_a=\frac{\pi_au'(C_a)}{u'(C_0)},\quad
 \phi_b=\frac{\pi_bu'(C_b)}{u'(C_0)},\quad
 \frac{\phi_a}{\phi_b}=\frac{\pi_au'(C_a)}{\pi_bu'(C_b)},\qquad
 1+r_F=\frac{u'(C_0)}{\pi_au'(C_a)+\pi_bu'(C_b)}.
 \]

 令 $D=\pi_au'(C_a)+\pi_bu'(C_b)$。由 $u'>0,u''<0,u'''<0$，得到
 \begin{center}
 \begin{tabular}{c|cccc}
 变化&$\phi_a$&$\phi_b$&$\phi_a/\phi_b$&$r_F$\\\hline
 $C_0$ 增加&上升&上升&不变&下降\\
 $C_a$ 增加&下降&不变&下降&上升\\
 $C_b$ 增加&不变&下降&上升&上升\\
 $(C_a,C_b)$ 同比增加&下降&下降&下降&上升\\
 等均值扩大两状态差距&下降&上升&下降&上升
 \end{tabular}
 \end{center}

 前三行直接来自边际效用递减。第四行中，对 $k>1$ 定义 $F(x)=u'(kx)/u'(x)$，则
 \[
 F'(x)=\frac{u'(kx)}{xu'(x)}[R(x)-R(kx)],\qquad
 R'(x)=-\frac{u''}{u'}-\frac{xu'''}{u'}+x\left(\frac{u''}{u'}\right)^2>0.
 \]

 因此 $C_a>C_b>0$ 时 $F(C_a)<F(C_b)$，同比增长后状态价格比下降。两未来状态的边际效用均下降，故 $D$ 下降、$r_F$ 上升。

 最后一行把“风险上升”明确为\textbf{保持期望总禀赋不变}的扩散：
 \[
 \pi_a\,dC_a+\pi_b\,dC_b=0,\qquad dC_a>0.
 \]

 沿这一变化方向，
 \[
 \frac{dD}{dC_a}=\pi_a[u''(C_a)-u''(C_b)]<0,
 \]

 因为 $u'''<0$ 使 $u''$ 严格递减。因此 $r_F$ 上升。只给出差距 $C_a-C_b$ 增加而不规定均值如何变化，尚未确定一种风险比较。
\end{enumerate}
\end{solution}

\par\noindent\begin{minipage}{\linewidth}
\begin{problem}
\WJTenExercise{exer:wj-ten-2}{**} 考虑一个含有参与者 $k=1,\dots,K$ 的 Arrow--Debreu 经济。参与者 $k$ 具有形式为 $u_{k,0}(c_{0})+\mathrm{E}[u_{k,1}(c_{1})]$ 的效用函数，对于所有 $k$，$u_{k,0}$ 和 $u_{k,1}$ 是严格递增和严格凹的。$f_{k}^{0}$ 和 $f_{k}$ 是在给定权重集合 $\lambda=\{\lambda_{1},\dots,\lambda_{K}\}$ 下对于参与者 $k$ 的 Pareto 最优分配规则。定义 $u_{0}$ 和 $u_{1}$ 如下
\[
    u_{0}'(x)=\lambda_{k}u_{k,0}'(f_{k}^{0}(x)),\qquad u_{1}'(x)=\lambda_{k}u_{k,1}'(f_{k}(x)).
\]
\begin{enumerate}[label=(\alph*)]
    \item 证明 $u_{0}$ 和 $u_{1}$ 是严格递增和严格凹的。
    \item 证明 $u_{0}(c_{0})+\mathrm{E}[u_{1}(c_{1})]$ 是代表性参与者的效用函数。
    \item 证明如果所有参与者都是绝对风险厌恶递减的，那么以上定义的代表性参与者的效用函数也是绝对风险厌恶递减的。
\end{enumerate}
\end{problem}
\end{minipage}\par
\Needspace{5\baselineskip}
\begin{solution}
% 原书练习10.2定义右侧漏写导数；题干已按边际效用校正，题号、其余设定保留。
对每一期，规划者条件把聚合边际效用与个人\textbf{边际效用}联系起来：
\[
 u_0'(x)=\lambda_ku_{k,0}'(f_k^0(x)),\qquad
 u_1'(x)=\lambda_ku_{k,1}'(f_k(x)).
\]

这里权重 $\lambda_k>0$ 固定，采用光滑内点模型，且 $u_k'>0,u_k''<0$。涉及 DARA 的导数时，效用三次可微。
\begin{enumerate}[label=(\alph*)]
 \item 由规划者条件，右侧对各人相同且为正，故 $u_0',u_1'>0$。分享函数严格递增；例如
 \[
 u_1''(x)=\lambda_ku_{k,1}''(f_k(x))f_k'(x)<0.
 \]

 当期效用也满足同样的导数关系，所以两期聚合效用均严格递增、严格凹。

 \item 令 $V_1(x)=\max_{\sum_kc_k=x}\sum_k\lambda_ku_{k,1}(c_k)$。包络定理给出
 \[
 V_1'(x)=\lambda_ku_{k,1}'(f_k(x))=u_1'(x),
 \]

 故 $u_1$ 与 $V_1$ 相差一个常数；$u_0$ 与当期规划最优值也相差一个常数。因此 $u_0(c_0)+\mathrm E[u_1(c_1)]$ 是代表性参与者的效用。

 \item 对任一期省略时间下标。记个人和代表性参与者的绝对风险容忍度为
 \[
 T_k(c)=-\frac{u_k'(c)}{u_k''(c)},\qquad
 T(x)=-\frac{u'(x)}{u''(x)}.
 \]

 对 $u'(x)=\lambda_ku_k'(f_k(x))$ 求导并相除，得 $f_k'=T_k(f_k)/T$。对各人求和，
 \[
 T(x)=\sum_kT_k(f_k(x)),\qquad
 T'(x)=\sum_kT_k'(f_k(x))f_k'(x)>0.
 \]

 DARA 等价于风险容忍度递增。每人的 $T_k'>0$，且 $f_k'>0$，故代表性参与者也有 $T'>0$，即 $A'(x)=-T'(x)/T(x)^2<0$。
\end{enumerate}
\end{solution}

\par\noindent\begin{minipage}{\linewidth}
\begin{problem}
\WJTenExercise{exer:wj-ten-3}{*} 在 1 期，经济有两个可能状态 $a$ 和 $b$。它们的发生的概率分别为 $\pi$ 和 $1-\pi$。市场有两只交易证券，证券 1 和 2，支付如下：
\begin{center}
\begin{tabular}{@{}c@{\hspace{2.5em}}c@{}}
证券 1：
\WJEventTree{1,1}
&
证券 2：
\WJEventTree{uC,dC}
\end{tabular}
\end{center}
其中 $C>0$ 且 $0<d<u$。因此，证券 1 是无风险债券而证券 2 是风险证券，我们可以称它为股票。经济中有参与者 $k=1,\dots,K$。参与者 $k$ 的禀赋是 $s_{k}C$ 单位的当期消费、0 单位的证券 1 和 $s_{k}$ 单位的证券 2，其中 $s_{k}>0$ 且 $\sum_{k}s_{k}=1$。所有参与者具有如下的期望效用函数：
\[
    U_{k}(c_{k})=\log c_{k,0}+\rho\bigl[\pi\log c_{k,1a}+(1-\pi)\log c_{k,1b}\bigr]
\]
其中 $\rho>0$。
\begin{enumerate}[label=(\alph*)]
    \item 证明证券市场是完全的。
    \item 求解每一参与者的最优消费/投资选择。
    \item 求解均衡证券价格。
    \item 将价格过程与期权定价中的二叉树模型相比较，你能得到什么结论？
\end{enumerate}
\end{problem}
\end{minipage}\par
\Needspace{5\baselineskip}
\begin{solution}
\begin{enumerate}[label=(\alph*)]
 \item 支付矩阵为 $X=\begin{bmatrix}1&uC\\1&dC\end{bmatrix}$，行列式为 $C(d-u)\ne0$，所以市场完全。

 \item 给定正状态价格，财富为 $w_k=s_kC(1+\phi_au+\phi_bd)$。最优条件为
 \[
 \frac1{c_{k,0}}=\lambda_k,\quad
 \frac{\rho\pi}{c_{k,a}}=\lambda_k\phi_a,\quad
 \frac{\rho(1-\pi)}{c_{k,b}}=\lambda_k\phi_b.
 \]

 代入预算得 $\lambda_k=(1+\rho)/w_k$，从而
 \[
 c_{k,0}=\frac{w_k}{1+\rho},\qquad
 c_{k,a}=\frac{\rho\pi w_k}{(1+\rho)\phi_a},\qquad
 c_{k,b}=\frac{\rho(1-\pi)w_k}{(1+\rho)\phi_b}.
 \]

 初始持仓为 $\bar\theta_k=(0,s_k)^{\mathsf T}$（按债券、股票排列）。由 $X(\bar\theta_k+\theta_k)=(c_{k,a},c_{k,b})^{\mathsf T}$ 解出净购买量：
 \[
 \theta_{k,B}=\frac{u c_{k,b}-d c_{k,a}}{u-d},\qquad
 \theta_{k,S}=\frac{c_{k,a}-c_{k,b}}{C(u-d)}-s_k.
 \]

 \item 令 $W=\sum_k w_k$ 为经济中的总财富。将各人的当期需求相加，当期出清要求
 \[
 C=\sum_k c_{k,0}=\frac{W}{1+\rho},
 \qquad W=(1+\rho)C.
 \]

 将两种状态下的个人需求分别相加，再代入 $W$，得到
 \[
 \begin{aligned}
 uC &= \sum_k c_{k,a}=\frac{\rho\pi W}{(1+\rho)\phi_a}
     =\frac{\rho\pi C}{\phi_a},\\
 dC &= \sum_k c_{k,b}=\frac{\rho(1-\pi)W}{(1+\rho)\phi_b}
     =\frac{\rho(1-\pi)C}{\phi_b}.
 \end{aligned}
 \]

 所以 $\phi_a=\rho\pi/u$、$\phi_b=\rho(1-\pi)/d$。按各状态的支付定价，
 \[
 B=\phi_a+\phi_b=\rho\left(\frac\pi u+\frac{1-\pi}d\right),
 \qquad S=\phi_a uC+\phi_b dC=\rho C.
 \]

 此时按禀赋计算的总财富为 $C+S=(1+\rho)C$，与当期出清求出的 $W$ 一致。

 代回需求，$c_{k,0}=s_kC$、$c_{k,a}=s_kuC$、$c_{k,b}=s_kdC$。因此 $\theta_{k,B}=\theta_{k,S}=0$，每人保留初始持仓。

 \item 二叉树期权定价通常给定债券、股票的\textbf{价格与支付}，再由复制关系确定状态价格。本题同时通过个人最优与市场出清内生求得证券价格，因此状态价格还反映时间偏好和实际概率。价格确定后，复制定价在两种分析中一致。
\end{enumerate}
\end{solution}

\par\noindent\begin{minipage}{\linewidth}
\begin{problem}
\WJTenExercise{exer:wj-ten-4}{**} 继续考虑练习~\ref{exer:wj-ten-3}（共同对数效用、两状态股票与债券经济）定义的经济，
\begin{enumerate}[label=(\alph*)]
    \item 构建代表性参与者。
    \item 计算在代表性参与者（即中央计划者）的效用函数中每个参与者的权重。它依赖于什么？为什么？
    \item 证明代表性参与者的偏好与禀赋在参与者之间的分布无关。
    \item 用基于消费的 CAPM 为两只证券定价。
    \item 计算无风险收益和股票的风险溢价。
    \item 利率和股票的风险溢价依赖于什么？解释它们对经济原生变量：当前消费水平、消费增长率的期望值 $\mathrm{E}[\tilde{C}_{1}/C_{0}]$、消费增长率的波动率 $\mathrm{Var}[\tilde{C}_{1}/C_{0}]$ 以及时间折现系数 $\rho$ 的依赖性？
\end{enumerate}
\end{problem}
\end{minipage}\par
\Needspace{5\baselineskip}
\begin{solution}
记 $\widetilde G=\widetilde C_1/C_0$，其取值为 $u,d$。本题的 $p$ 表示股票\textbf{净收益率}的风险溢价。
\begin{enumerate}[label=(\alph*)]
 \item 对数效用的最优分享为 $c_k=\mu_kC/\sum_j\mu_j$。取 $\sum_k\mu_k=1$，由均衡消费 $c_k=s_kC$ 得 $\mu_k=s_k$。聚合效用为
 \[
 \log C_0+\rho\bigl[\pi\log C_{1a}+(1-\pi)\log C_{1b}\bigr]
 \]
 加上一个与消费无关的常数。

 \item 规划权重为 $\mu_k=s_k$，由个人禀赋份额决定。较高权重对应较高的最优消费份额。

 \item 禀赋份额只改变加总中的常数项，代表性参与者的消费排序与 $s_k$ 无关。

 \item 随机贴现因子为 $\widetilde m=\rho/\widetilde G$。由\ifdefined\WJHandoutKnowledge 式~\eqref{eq:wj-10-16}\else 证券定价公式\fi（定价核加权的证券支付定价公式），
 \[
 B=\rho\mathrm E[1/\widetilde G],\qquad
 S=\mathrm E[(\rho/\widetilde G)C\widetilde G]=\rho C.
 \]

 \item 债券与股票的期望总收益率分别为
 \[
 1+r_F=\frac1{\rho\mathrm E[1/\widetilde G]},\qquad
 1+\bar r_S=\frac{\mathrm E[\widetilde G]}\rho.
 \]

 因此
 \[
 \begin{aligned}
 p=\bar r_S-r_F
 &=\frac1\rho\left(\mathrm E[\widetilde G]-\frac1{\mathrm E[1/\widetilde G]}\right)\\
 &=\frac{\pi(1-\pi)(u-d)^2}{\rho[\pi d+(1-\pi)u]}>0.
 \end{aligned}
 \]

 这是算术均值与调和均值之差；$u=d$ 时风险溢价为零。

 \item 给定增长率分布，消费水平 $C$ 按比例改变股票价格与支付，收益率和风险溢价保持不变。增大 $\rho$ 使 $r_F$ 下降，并按 $1/\rho$ 的比例降低风险溢价。

 固定状态概率及价差 $u-d$，把两个增长率一起上移，$\mathrm E[1/\widetilde G]$ 下降，故 $r_F$ 上升；上式分母增大，风险溢价下降。固定均值而扩大价差，则构成均值保持扩散。函数 $1/G$ 严格凸，故债券价格上升、无风险利率下降；股票期望收益率保持不变，风险溢价上升。

 令 $\tilde g=\widetilde G-1$。在其取值较小时，二阶展开给出
 \[
 1+r_F\approx\frac1{\rho[1-\bar g+\operatorname{Var}(\tilde g)+\bar g^2]},\qquad
 p\approx\frac{\operatorname{Var}(\tilde g)}\rho.
 \]

 上式给出债券的总收益率。本模型始终有 $S=\rho C$；保持 $C,\rho$ 不变而增加未来风险时，股票现价保持不变，收益率溢价的上升来自无风险利率下降。
\end{enumerate}
\end{solution}

\par\noindent\begin{minipage}{\linewidth}
\begin{problem}
\WJTenExercise{exer:wj-ten-5}{} 继续前面练习中的经济。
\begin{enumerate}[label=(\alph*)]
    \item 求解以证券 2 为标的资产、执行价格为 $K$ 的看涨期权的价格。
    \item 证明基于消费的 CAPM 也适用于期权。
\end{enumerate}
\end{problem}
\end{minipage}\par
\Needspace{5\baselineskip}
\begin{solution}
\begin{enumerate}[label=(\alph*)]
 \item 记 $H_a=(uC-K)_+$、$H_b=(dC-K)_+$。复制期权的组合中，债券、股票的数量分别为
 \[
 \theta_B=\frac{uH_b-dH_a}{u-d},\qquad
 \theta_S=\frac{H_a-H_b}{C(u-d)}.
 \]

 所以 $V=\theta_BB+\theta_SS=\phi_aH_a+\phi_bH_b$。按执行价写成
 \[
 V=\begin{cases}
 \rho C-K\rho(\pi/u+(1-\pi)/d),&K\le dC,\\
 (\rho\pi/u)(uC-K),&dC<K<uC,\\
 0,&K\ge uC.
 \end{cases}
 \]

 \item 用 $\widetilde m=\rho C/\widetilde C_1$ 定价，
 \[
 V=\mathrm E[\widetilde m(\widetilde C_1-K)_+]
 =\frac{\rho\pi}{u}(uC-K)_+
  +\frac{\rho(1-\pi)}d(dC-K)_+.
 \]

 这与复制价格对任意执行价都相同，因此消费定价公式同样适用于期权。
\end{enumerate}
\end{solution}

\par\noindent\begin{minipage}{\linewidth}
\begin{problem}
\WJTenExercise{exer:wj-ten-6}{**} 在 1 期，经济有两个可能状态 $a$ 和 $b$。它们的发生概率分别为 $\pi$ 和 $1-\pi$。市场有两只交易证券，证券 1 和 2，支付如下：
\begin{center}
\begin{tabular}{@{}c@{\hspace{2.5em}}c@{}}
证券 1：
\WJEventTree{1,1}
&
证券 2：
\WJEventTree{C_{1a},C_{1b}}
\end{tabular}
\end{center}
% 原书沿用前题C、u、d，未在本题定义；按本题支付补足正消费及完全性条件。
其中 $C_0>0$ 且 $0<C_{1b}<C_{1a}$。因此，证券 1 是无风险债券而证券 2 是风险证券，我们可以称它为股票。经济中有参与者 $k=1,\dots,K$。参与者 $k$ 的禀赋是 $s_{k}C_{0}$ 单位的当期消费、0 单位的证券 1 和 $s_{k}$ 单位的证券 2，其中 $s_{k}>0$ 且 $\sum_{k}s_{k}=1$。所有参与者具有如下的期望效用函数：
\[
    U_{k}(c_{k})=-e^{-\alpha c_{k,0}}-\rho\bigl[\pi e^{-\alpha c_{k,1a}}+(1-\pi)e^{-\alpha c_{k,1b}}\bigr]
\]
其中 $\rho>0$、$\alpha>0$、$0<\pi<1$。本题按消费非负约束不绑定的内点模型求解。
\begin{enumerate}[label=(\alph*)]
    \item 证明证券市场是完全的。
    \item 求解每一参与者的最优消费/投资选择。
    \item 求解均衡证券价格。
\end{enumerate}
\end{problem}
\end{minipage}\par
\Needspace{5\baselineskip}
\begin{solution}
% 原书在本题沿用了前题的C、u、d文字，却以C_1a、C_1b给出支付；这里明确实际使用的完全性条件。
两状态股票支付不同，债券和股票可以复制任意状态支付。以下利用题设内点条件求解个人选择与均衡。
\begin{enumerate}[label=(\alph*)]
 \item $X=\begin{bmatrix}1&C_{1a}\\1&C_{1b}\end{bmatrix}$ 的行列式非零，市场完全。

 \item 给定状态价格，记
 \[
 B=\phi_a+\phi_b,\quad S=\phi_aC_{1a}+\phi_bC_{1b},\quad
 W=C_0+S,\quad w_k=s_kW.
 \]

 个人最优条件为
 \[
 \alpha e^{-\alpha c_{k,0}}=\lambda_k,\quad
 \rho\pi\alpha e^{-\alpha c_{k,a}}=\lambda_k\phi_a,\quad
 \rho(1-\pi)\alpha e^{-\alpha c_{k,b}}=\lambda_k\phi_b.
 \]

 令 $d_a=\alpha^{-1}\log(\rho\pi/\phi_a)$、$d_b=\alpha^{-1}\log(\rho(1-\pi)/\phi_b)$。前后相除得 $c_{k,a}=c_{k,0}+d_a$、$c_{k,b}=c_{k,0}+d_b$。代入预算，
 \[
 c_{k,0}=\frac{s_kW-\phi_ad_a-\phi_bd_b}{1+B},\quad
 c_{k,a}=c_{k,0}+d_a,\quad c_{k,b}=c_{k,0}+d_b.
 \]

 初始持仓为 $\bar\theta_k=(0,s_k)^{\mathsf T}$。由 $X(\bar\theta_k+\theta_k)=(c_{k,a},c_{k,b})^{\mathsf T}$，得到债券和股票的净购买量：
 \[
 \theta_{k,B}=\frac{C_{1a}c_{k,b}-C_{1b}c_{k,a}}{C_{1a}-C_{1b}},\qquad
 \theta_{k,S}=\frac{c_{k,a}-c_{k,b}}{C_{1a}-C_{1b}}-s_k.
 \]

 \item 将消费关系对各人求和，出清给出 $Kd_a=C_{1a}-C_0$、$Kd_b=C_{1b}-C_0$。因此
 \[
 \phi_a=\rho\pi e^{-\frac\alpha K(C_{1a}-C_0)},\qquad
 \phi_b=\rho(1-\pi)e^{-\frac\alpha K(C_{1b}-C_0)}.
 \]

 由此计算 $B,S,W$。令 $\delta_k=(s_k-1/K)W/(1+B)$，均衡消费可简写为
 \[
 c_{k,0}=C_0/K+\delta_k,\quad
 c_{k,a}=C_{1a}/K+\delta_k,\quad
 c_{k,b}=C_{1b}/K+\delta_k.
 \]

 代回净购买量，得到
 \[
 \theta_{k,B}=\delta_k,\qquad \theta_{k,S}=1/K-s_k.
 \]

 因 $\sum_k\delta_k=0$，债券总持仓为零，股票总持仓为1，净交易总量均为零。上述消费若均严格为正，就给出非负消费模型的内点均衡。
\end{enumerate}
\end{solution}

\par\noindent\begin{minipage}{\linewidth}
\begin{problem}
\WJTenExercise{exer:wj-ten-7}{**} 继续考虑练习~\ref{exer:wj-ten-6}（共同指数效用、两状态股票与债券的内点模型）定义的经济
\begin{enumerate}[label=(\alph*)]
    \item 构建代表性参与者。
    \item 计算在代表性参与者（即中央计划者）的效用函数中每个参与者的权重。它依赖于什么？为什么？
    \item 证明代表性参与者的偏好与禀赋在参与者之间的分布无关。
    \item 用基于消费的 CAPM 为两只证券定价。
    \item 计算无风险收益和股票的风险溢价，现在风险溢价的定义为 $\mathrm{E}[\tilde{S}_{1}]-(1+r_{F})S_{0}$。即一股股票的绝对超额收益。
    \item 利率和股票的风险溢价依赖于什么？解释它们对经济原生变量：当前消费水平、消费增长的期望值 $\mathrm{E}[\tilde{C}_{1}/C_{0}]$、消费增长的波动率 $\mathrm{Var}[\tilde{C}_{1}/C_{0}]$ 以及时间折现系数 $\rho$ 的依赖性？
    \item 求解以证券 2 为标的资产、执行价格为 $L$ 的看涨期权的价格。
\end{enumerate}
\end{problem}
\end{minipage}\par
\Needspace{5\baselineskip}
\begin{solution}
沿用练习~\ref{exer:wj-ten-6} 的内点模型，令 $a=\alpha/K$、$B=\phi_a+\phi_b$、$W=C_0+S_0$，并记 $\delta_k=(s_k-1/K)W/(1+B)$。
\begin{enumerate}[label=(\alph*)]
 \item 均衡中 $c_{k,s}=C_s/K+\delta_k$。规划者边际效用满足
 \[
 \mu_k\alpha e^{-\alpha c_{k,s}}
 =\alpha\mu_ke^{-\alpha\delta_k}e^{-aC_s}.
 \]

 取 $\mu_k\propto e^{\alpha\delta_k}$，各人的加权边际效用相等，聚合效用经正比例标准化为
 \[
 -e^{-aC_0}-\rho[\pi e^{-aC_{1a}}+(1-\pi)e^{-aC_{1b}}].
 \]

 \item 归一化后的权重可写为
 \[
 \mu_k=\frac{\exp(\alpha Ws_k/(1+B))}
 {\sum_j\exp(\alpha Ws_j/(1+B))}.
 \]

 它随禀赋份额 $s_k$ 上升；正指数中的 $W$ 反映财富单位，比较的对象是个人财富差异。

 \item 财富份额影响聚合效用的正比例系数，代表性偏好的形状由 $a=\alpha/K$ 决定。因此在本题共同时间折现下，定价与 $s_k$ 无关。

 \item 定价核为 $\widetilde m=\rho e^{-a(\widetilde C_1-C_0)}$。由\ifdefined\WJHandoutKnowledge 式~\eqref{eq:wj-10-16}\else 证券定价公式\fi（定价核加权的证券支付定价公式），
 \[
 B=\rho\mathrm E[e^{-a(\widetilde C_1-C_0)}],\qquad
 S_0=\rho\mathrm E[e^{-a(\widetilde C_1-C_0)}\widetilde C_1].
 \]

 \item 无风险\textbf{净}利率为 $r_F=1/B-1$。本题风险溢价为一股股票的\textbf{绝对超额收益}。记 $\Delta=C_{1a}-C_{1b}>0$，则
 \[
 \begin{aligned}
 p&=\pi C_{1a}+(1-\pi)C_{1b}-S_0/B\\
 &=\left(\pi-\frac{\phi_a}{B}\right)\Delta\\
 &=\frac{\pi(1-\pi)(e^{a\Delta}-1)}{\pi+(1-\pi)e^{a\Delta}}\,\Delta>0.
 \end{aligned}
 \]

 它等于实际概率与风险中性概率之差乘以状态支付差。

 \item 给定未来\textbf{增长率分布}，令 $\tilde g=\widetilde C_1/C_0-1$、$\ell=aC_0$。当 $|\ell\tilde g|$ 一致较小时，
 \[
 B\approx\rho[1-\ell\bar g+\tfrac12\ell^2(\operatorname{Var}(\tilde g)+\bar g^2)],
 \quad
 r_F\approx\frac1{\rho[1-\ell\bar g+\tfrac12\ell^2(\operatorname{Var}(\tilde g)+\bar g^2)]}-1.
 \]

 展开变量为 $\alpha C_0\tilde g/K$，人均消费增大也会放大边际效用的变化。

 $\rho$ 增大使债券价格上升、无风险利率下降；绝对风险溢价中 $\rho$ 相消。固定未来支付差而将两个未来消费水平一起上移，债券价格下降、利率上升，绝对风险溢价保持不变。固定概率和平均未来消费而扩大 $\Delta$，指数函数的凸性使债券价格上升、利率下降，上式的风险溢价则严格增加。

 给定增长率分布改变 $C_0$ 时，利率方向由
 \[
 \frac{\partial r_F}{\partial C_0}
 =\frac{\rho a}{B^2}\mathrm E[\tilde g e^{-aC_0\tilde g}]
 \]

 决定。与此同时，$\Delta=C_0(g_a-g_b)$ 随消费规模增大，绝对风险溢价也增加。若改为保持未来消费\textbf{水平}不变而提高 $C_0$，则债券价格上升、利率下降，绝对风险溢价不变。这两种比较分别固定增长率与未来水平。

 \item % 原书本小题执行价也记K，与参与者人数冲突；题干、解答均局部记为L。
执行价为 $L$ 的看涨期权价格为
 \[
 V(L)=\phi_a(C_{1a}-L)_++\phi_b(C_{1b}-L)_+
 =\rho\mathrm E[e^{-a(\widetilde C_1-C_0)}(\widetilde C_1-L)_+].
 \]

 它与由债券、股票复制该支付的成本相同，适用于全部执行价。
\end{enumerate}
\end{solution}

